\section*{Chapitre \ref{ch:b2:chromatography} --- Chromatographie}

\begin{solution}{exo:b2:chromatography:1}
$k = (5.0 - 1.0)/1.0 = 4.0$~; fraction du temps passée dans la \omterm{def:b2:chromatography:principle}{phase mobile}~: $1/(1 + k) = 0.20$.
\end{solution}

\begin{solution}{exo:b2:chromatography:2}
$N = 16(6.0/0.30)^2 = 16 \times 400 = 6400$.
\end{solution}

\begin{solution}{exo:b2:chromatography:3}
$H = \qty{250}{mm}/10\,000 = \qty{0.025}{mm} = \qty{25}{\micro m}$.
\end{solution}

\begin{solution}{exo:b2:chromatography:4}
Éthanol dans le sang~: CPG (volatil). Protéine~: CLHP. Caféine dans une boisson~: CLHP (dans l'eau, pas
assez volatile sans préparation). Benzène dans l'essence~: CPG. Sucres~: CLHP (ils se décomposent avant
de bouillir).
\end{solution}

\begin{solution}{exo:b2:chromatography:5}
$R_s = 2(4.6 - 4.0)/(0.40 + 0.50) = 1.2/0.90 \approx 1.3$~: pas tout à fait séparés à la ligne de base (au-dessous de 1.5).
\end{solution}

\begin{solution}{exo:b2:chromatography:6}
$\sqrt N = 4R_s\,\dfrac{\alpha}{\alpha - 1}\,\dfrac{1 + k_2}{k_2} = 4 \times 1.5 \times 11 \times 1.25 =
82.5$, d'où $N \approx 6.8 \times 10^3$.
\end{solution}

\begin{solution}{exo:b2:chromatography:7}
$u_{\mathrm{opt}} = \sqrt{8/2} = \qty{2}{mm/s}$~; $H_{\min} = 5 + 2\sqrt{16} = \qty{13}{\micro m}$.
\end{solution}

\begin{solution}{exo:b2:chromatography:8}
Le phénol (\ce{OH} polaire, liaisons hydrogène avec l'eau) d'abord, puis le benzène, puis le toluène (un
\ce{CH3} de plus, plus apolaire). Davantage de méthanol rend la \omterm{def:b2:chromatography:principle}{phase mobile} moins polaire~: les trois
composés sortent plus tôt, dans le même ordre.
\end{solution}

\begin{solution}{exo:b2:chromatography:9}
$F = (860/400)/(20.0/10.0) = 2.15/2.00 = 1.075$. Échantillon~: $c_X = (645/410) \times 10.0/1.075 \approx
\qty{14.6}{mg/L}$.
\end{solution}

\begin{solution}{exo:b2:chromatography:10}
$R_s \propto \sqrt N \propto \sqrt L$~: il faut multiplier $L$ par $(1.5/1.06)^2 \approx 2.0$, ce qui double
la durée de l'analyse et la pression. Autres leviers~: la sélectivité (\omterm{def:b2:chromatography:principle}{phase mobile}, \omterm{def:b2:chromatography:principle}{phase stationnaire},
température), des particules plus fines, le débit optimal.
\end{solution}

\begin{solution}{exo:b2:chromatography:11}
Les pics sont distants de $4\sigma R_s$, si bien que le point médian est à $2\sigma R_s$ de chacun. $R_s = 1$~: chaque pic
apporte $\mathrm e^{-(2)^2/2} = \mathrm e^{-2} = 0.135$~; la vallée est à $0.27$ de la hauteur des pics.
$R_s = 1.5$~: $2\mathrm e^{-(3)^2/2} = 2\mathrm e^{-4.5} \approx 0.022$, environ 2~\%~: retour à la ligne de
base en pratique.
\end{solution}

\begin{solution}{exo:b2:chromatography:12}
$k_2/(1+k_2)$~: 0.67, 0.83, 0.91~; durée en unités de $t_M$~: 3, 6, 11. Passer de 2 à 5 gagne 25~\% de
\omterm{def:b2:chromatography:separation}{résolution} pour une durée doublée~; de 5 à 10, 9~\% pour une durée encore presque doublée. Un $k$ de 2 à 5
environ est le compromis habituel.
\end{solution}

\begin{solution}{pb:b2:chromatography:1}
\textbf{1.} Stationnaire~: grains de silice greffés de longues chaînes alkyle, apolaires. Mobile~: eau et
méthanol, polaire.
\textbf{2.} Ils sont très polaires et restent dans la \omterm{def:b2:chromatography:principle}{phase mobile}~: $k \approx 0$.
\textbf{3.} La caféine porte un groupe méthyle de plus que la théophylline, et la théophylline une liaison
\ce{N-H} qui forme des liaisons hydrogène avec l'eau~: la théophylline est plus polaire et sort la première.
\textbf{4.} Elle a une structure proche (réponse et comportement semblables), est absente de la boisson, et
sort près de la caféine tout en en étant résolue.
\textbf{5.} Tous deux diminueraient~: une \omterm{def:b2:chromatography:principle}{phase mobile} moins polaire concurrence mieux la \omterm{def:b2:chromatography:principle}{phase stationnaire}.
\textbf{6.} Le temps que met la \omterm{def:b2:chromatography:principle}{phase mobile} à traverser la colonne, c'est-à-dire le temps pendant lequel un
composé se déplace.
\textbf{7.} $k_{\mathrm{theo}} = (3.0 - 1.0)/1.0 = 2.0$~; $k_{\mathrm{caf}} = 2.4$.
\textbf{8.} $\alpha = 2.4/2.0 = 1.2$.
\textbf{9.} Caféine~: $N = 16(3.4/0.20)^2 = 4624 \approx 4.6 \times 10^3$~; théophylline~: $16(3.0/0.18)^2
\approx 4.4 \times 10^3$.
\textbf{10.} $H = \qty{150}{mm}/4624 \approx \qty{0.032}{mm} = \qty{32}{\micro m}$.
\textbf{11.} $R_s = 2(3.4 - 3.0)/(0.18 + 0.20) = 0.80/0.38 \approx 2.1$.
\textbf{12.} $\frac{\sqrt{4624}}{4} \times \frac{0.2}{1.2} \times \frac{2.4}{3.4} = 17 \times 0.167 \times
0.706 \approx 2.0$~; le petit écart vient de l'hypothèse de largeurs égales faite dans l'équation.
\textbf{13.} Oui~: $R_s > 1.5$.
\textbf{14.} $N$ est divisé par deux~: $R_s \approx 2.1/\sqrt2 \approx 1.5$, juste à la séparation à la ligne de base.
\textbf{15.} Divisée par deux~: la caféine sort à \qty{1.7}{min}.
\textbf{16.} $R_s \approx 2.1\sqrt2 \approx 3.0$, pour une durée et une pression doublées~: un gaspillage ici.
\textbf{17.} La composition de la \omterm{def:b2:chromatography:principle}{phase mobile} (fraction de méthanol, pH, autre solvant organique) ou la
\omterm{def:b2:chromatography:principle}{phase stationnaire}.
\textbf{18.} Peu~: $k_2/(1 + k_2) = 0.71$ déjà~; pour $k_2 = 5$, il vaudrait 0.83 (+18~\%), pour une
durée d'analyse multipliée par $6/3.4 \approx 1.8$.
\textbf{19.} $F = (1500/1000)/(50.0/40.0) = 1.50/1.25 = 1.20$.
\textbf{20.} Les deux pics proviennent de la même injection~: le volume se simplifie dans le rapport des aires.
\textbf{21.} $c = (970/1010) \times 40.0/1.20 \approx \qty{32.0}{mg/L}$.
\textbf{22.} Dilution au dixième~: \qty{320}{mg/L}.
\textbf{23.} $A_{IS}$ serait trop grande, et le résultat pour la caféine trop faible.
\textbf{24.} Une série d'étalons de caféine injectés exactement dans les mêmes conditions que l'échantillon,
une droite d'étalonnage de l'aire en fonction de la concentration, et des volumes d'injection reproductibles.
\textbf{25.} $\qty{320}{mg/L} \times \qty{0.250}{L}$~: $\boldsymbol{\approx \qty{80}{mg}}$ de caféine par canette
(données de l'exercice).
\end{solution}
